% !TEX root = ../main.tex
\begin{center}
    \large
    \begin{singlespace}    
        \textbf{\chineseTitle{}} \\[0.5cm]
    \end{singlespace}

    \begin{singlespace}    
    學生：\studentCnName  \hspace{2.5cm}  指導教授：\advisorCnName \hspace{0.1cm} 博士 \\

    \ifdefined\advisorCnNameB % 如果有共同指導教授
        \hspace{9.6cm}  \advisorCnNameB ~ 博士 \\
    \fi
    \end{singlespace}

    \vspace{0.5cm}

    \begin{singlespace}
        國立陽明交通大學\\
        \NameofDepartmentInstituteCN\\[0.2cm]
    \end{singlespace}

    \textbf{摘\hspace{1cm}要} \\[0.5cm]

\end{center}

\normalsize 

% TODO 最後寫（等第 7 章有數字之後）。結構建議照 Introduction 的四段：
%      情境 → 現行兩條路都不對 → 把推論從執行期移到編譯期 → 結果數字。
%      注意適用條件要出現在摘要裡：重複性模板實例化任務、輸出可由結構化來源資料
%      與明確操作規則確定性推導。
（中文摘要待撰寫。）

\vspace{1cm}

% 中文摘要及關鍵詞 5-7 個
% TODO 候選：模板實例化、可驗證規格、確定性執行、大型語言模型、文件自動化、程式合成
\textbf{關鍵字：}（待定，5--7 個）
